% Einheit 4 – Gleichung bestimmen (Nr. 34–44)
\setcounter{aufgabe}{33}
\einheitenkopf[e4][4 Gleichung bestimmen]{Einheit 4 von 5 · Gleichung bestimmen}
\zweigzeile{Hier lernst du, die Gleichung einer Geraden aus dem Graphen, aus Steigung und Punkt oder aus zwei Punkten zu bestimmen und Schnittpunkte zu berechnen · neu in diesem Jahr · P10 · baut auf: $m$ und $n$ ablesen (Einheit 2), Punktprobe (Einheit 3), Gleichungen lösen, Terme zusammenfassen}

\verfahren{Gleichung aus dem Graphen}

\begin{aufgabe}{Ich kann die Gleichung einer Geraden aus dem Graphen bestimmen.\newline Lies $m$ und $n$ ab und gib die Gleichung an.}

\begin{ksys}[xmin=-4,xmax=4,ymin=-5,ymax=6,ablesen]
\gerade[-2]{1}{3}{r}
\gerade[2.5]{3}{-3}{s}
\gerade[3.5]{2}{-4}{t}
\end{ksys}

\begin{teile}
\teil Gerade $r$: \quad \feld{m} \quad \feld{n} \quad \feld{y}
\teil Gerade $s$: \quad \feld{m} \quad \feld{n} \quad \feld{y}
\teil Gerade $t$: \quad \feld{m} \quad \feld{n} \quad \feld{y}
\end{teile}
\end{aufgabe}

\verfahren{Gleichung aus Steigung und Punkt}

\begin{aufgabe}{Ich kann die Gleichung aus der Steigung und einem Punkt bestimmen.\newline Die Gerade hat die Steigung $m$ und geht durch den Punkt $P$. Bestimme $n$ und gib die Gleichung an.}
\rechenplatz{3}
\begin{gleichungsraster}[3]
\gl{m = 2;\ \ P(3|8)} & \gl{m = 3;\ \ P(2|5)} \\
\gl{m = -1;\ \ P(4|1)} & \gl{m = 0{,}5;\ \ P({-4}|1)} \\
\end{gleichungsraster}
\end{aufgabe}

\verfahren{Gleichung aus zwei Punkten}

\begin{aufgabe}{Ich kann die Gleichung aus zwei Punkten bestimmen.\newline Die Gerade geht durch $A$ und $B$. Bestimme $m$ und $n$ und gib die Gleichung an.}
\rechenplatz{4}
\begin{gleichungsraster}[3]
\gl{A(0|4),\ \ B(2|10)} & \gl{A(0|{-2}),\ \ B(3|4)} \\
\gl{A(1|5),\ \ B(4|11)} & \\
\end{gleichungsraster}
\end{aufgabe}

\begin{aufgabe}{Ich kann die Gleichung aus zwei Punkten bestimmen – weiter.}
\begin{gleichungsraster}[3]
\gl{A(2|7),\ \ B(5|1)} & \gl{A({-4}|1),\ \ B(2|4) \quad \text{(P10 2025 OS)}} \\
\end{gleichungsraster}
\end{aufgabe}

\begin{aufgabe}{Ich kann eine Gerade durch zwei Punkte zeichnen.\newline Trage die Punkte ein und zeichne die Gerade durch sie.}
\begin{teile}
\teil $P(0|{-1})$ und $Q(2|3)$
\teil $P({-2}|4)$ und $Q(2|2)$
\teil $C({-3}|5)$ und $D(3|1)$. Bestimme danach $m$ und $n$ und gib die Gleichung an. (P10 2025 OS) \quad \feld{y}
\end{teile}

\begin{ksys}[xmin=-4,xmax=4,ymin=-3,ymax=6]
\end{ksys}

\end{aufgabe}

\verfahren{Schnittpunkt zweier Geraden}

\begin{aufgabe}{Ich kann den Schnittpunkt zweier Geraden berechnen.\newline Berechne den Schnittpunkt der beiden Geraden.}
\rechenplatz{4}
\begin{gleichungsraster}[3]
\gl{y = x + 1;\ \ y = -x + 5} & \gl{y = 2x - 1;\ \ y = x + 3} \\
\gl{y = 3x + 7;\ \ y = -2x - 3} & \\
\end{gleichungsraster}
\end{aufgabe}

\begin{aufgabe}{Ich kann den Schnittpunkt zweier Geraden berechnen – weiter.}
\begin{gleichungsraster}[3]
\gl{y = 0{,}5x + 1;\ \ y = -x + 4} & \gl{y = 4x - 3;\ \ y = 4x + 1} \\
\end{gleichungsraster}
\end{aufgabe}

\verfahren{Prüfen, begründen, anwenden}

\begin{aufgabe}{Ich finde den Fehler beim Bestimmen einer Geradengleichung.\newline Finde den Fehler, benenne ihn und korrigiere die Rechnung.}
\begin{teile}
\teil Ali bestimmt die Steigung der Geraden durch $P(2|5)$ und $Q(6|13)$:
\rechnung{m &= \frac{13 - 5}{2 - 6} = \frac{8}{-4} = -2}
\teil Mehmet bestimmt die Gerade mit $m = 3$ durch $P(2|4)$:
\rechnung{2 &= 3 \cdot 4 + n \\ n &= -10 \\ y &= 3x - 10}
\end{teile}
\end{aufgabe}

\begin{aufgabe}{Ich kann begründen, wie zwei Geraden zueinander liegen.\newline Entscheide ohne Rechnung und begründe.}
\begin{teile}
\teil Haben $y = 3x + 6$ und $y = 3x - 5$ einen Schnittpunkt?
\teil Wo schneiden sich $y = 2x + 4$ und $y = -0{,}5x + 4$?
\end{teile}
\end{aufgabe}

\begin{aufgabe}{Ich kann eine Geradengleichung aus einer Sachsituation bestimmen.\newline Ein Heißluftballon sinkt gleichmäßig. Nach 2 Minuten ist er 840\,m hoch, nach 6 Minuten 600\,m. $x$ ist die Zeit in Minuten, $y$ die Höhe in Metern.}
\begin{teile}
\teil Bestimme die Gleichung. (P10 2015 OS) \quad $y = $ \leerfeld
\teil Wie hoch war der Ballon, als er zu sinken begann? \quad \leerfeld[m]
\teil Nach wie vielen Minuten landet er? \quad \leerfeld[min]
\end{teile}
\end{aufgabe}

\begin{aufgabe}{Ich kann mit einem Schnittpunkt eine Sachaufgabe lösen.\newline Anna fährt mit dem Rad los, 20\,km in jeder Stunde. Ben ist zur selben Zeit schon 12\,km weiter auf derselben Strecke und fährt 12\,km in jeder Stunde in dieselbe Richtung. $x$ ist die Zeit in Stunden, $y$ der Weg ab Annas Start in km.}
\begin{teile}
\teil Stelle für Anna und für Ben die Gleichung auf.
\teil Wann und wo holt Anna Ben ein?
\end{teile}
\end{aufgabe}
